% from ega0-19-fr.tex
\oldpage[0\textsubscript{IV-1}]{81}

\begin{proof}
(i) résulte aussitôt de ce que $f^{-1}(0)$ est un voisinage de $0$ dans le
$A$-module produit $A^{\mathbb{N}^{(I)}}$, d'où $f(T^\lambda)=0$ sauf pour un
nombre fini de valeurs de $\lambda\in\mathbb{N}^{(I)}$. Pour prouver (ii), il
suffit de remarquer que l'anneau de polynômes
$B'=A[T_\alpha]_{\alpha\in I}$ est dense dans $B$~; l'existence et
l'unicité de la restriction $g|B'$ sont triviales et sa continuité résulte de
l'hypothèse que les $f(T_\alpha)$ pour $\alpha\in H$ sont nilpotents,
car si $(f(T_\alpha))^n=0$ pour tout $\alpha\in H$, on a
$g(T^\lambda)=0$ pour tout $\lambda=(\lambda_\alpha)_{\alpha\in I}$ de
support fini sauf ceux tels que $\lambda_\alpha=0$ pour
$\alpha\notin H$ et $\lambda_\alpha\leq n$ pour
$\alpha\in H$, c'est-à-dire sauf pour un nombre fini de valeurs de
$\lambda$.
\end{proof}

\begin{proof}[Démonstration de la proposition~{\rm(19.3.4)}]
Ce lemme étant établi, et les notations $C$ et $\mathfrak J$ ayant le même
sens que dans (19.3.1), on a $u(T_\alpha)=0$ sauf pour
$\alpha\in H$, $H$ étant une partie finie de $I$, et les
$z_\alpha=u(T_\alpha)$ pour $\alpha\in H$ sont nilpotents dans
$C/\mathfrak J$~; comme $\mathfrak J$ est nilpotent, il existe une famille
$(x_\alpha)_{\alpha\in H}$ d'éléments nilpotents de $C$ dont les images
canoniques dans $C/\mathfrak J$ sont les $z_\alpha$~; si $v$ est le
$A$-homomorphisme continu de $B$ dans $C$ tel que
$v(T_\alpha)=0$ pour $\alpha\notin H$, $v(T_\alpha)=x_\alpha$ pour
$\alpha\in H$, il est clair que $u$ se factorise en
$B\xrightarrow{v}C\to C/\mathfrak J$.
\end{proof}

\begin{proposition}[19.3.5]
\label{0.19.3.5-fr}
Soit $A$ un anneau topologique.
\begin{enumerate}
  \item[(i)] $A$ est une $A$-algèbre formellement lisse.
  \item[(ii)] Si $B$ est une $A$-algèbre formellement lisse et $C$ une
  $B$-algèbre formellement lisse, alors $C$ est une $A$-algèbre formellement
  lisse.
  \item[(iii)] Soient $B$ une $A$-algèbre formellement lisse, $A'$ une
  $A$-algèbre topologique~; alors la $A'$-algèbre topologique
  $B\otimes_A A'$ ($\mathbf 0_{\rm I}$, 7.7.5 et 7.7.6) est formellement
  lisse.
  \item[(iv)] Soient $B$ une $A$-algèbre topologique, $S$ (resp. $T$) une
  partie multiplicative de $A$ (resp. $B$) telle que l'image canonique de $S$
  dans $B$ soit contenue dans $T$. Si $B$ est une $A$-algèbre formellement
  lisse, alors $T^{-1}B$ est une $S^{-1}A$-algèbre formellement lisse.
  \item[(v)] Soient $B_i$ ($1\leq i\leq n$) des $A$-algèbres topologiques.
  Pour que $\prod_{i=1}^n B_i$ soit une $A$-algèbre formellement lisse, il
  faut et il suffit que chacune des $B_i$ le soit.
\end{enumerate}
\end{proposition}

